% Draft for subsection 3.9.2; insert after cmp:subsec:ch3-policy-geometry-map.
% All proofs are deliberately retained here pending editorial review.
\subsection{سیاست‌های هم‌مرکز}
\label{cmp:subsec:ch3-geometry-parametric-fibers}

در بخش پیشین، محدودیت‌های دست‌یابی به آرایش‌های دلخواه بازنمایی‌ها و مراکز طبقه‌ها بررسی شدند. اکنون مسئله را معکوس می‌کنیم: اگر مراکز از یک سیاست مجاز حاصل شده باشند و در نتیجه دست‌یافتنی بودن آن‌ها معلوم باشد، چه سیاست‌های دیگری می‌توانند همان مراکز را بازتولید کنند؟ آیا این سیاست‌ها الزاماً بازنمایی‌های یکسانی دارند یا می‌توان با حفظ مراکز، آرایش نمونه‌ها را تغییر داد؟

اهمیت این پرسش در آن است که ثابت بودن مراکز، معیارهای وابسته صرفاً به مراکز و، در قاعدهٔ نزدیک‌ترین مرکز، مجموعه‌های جداکننده را ثابت نگه می‌دارد؛ اما لزوماً موقعیت بازنمایی‌های نمونه‌ها، تراکم درون‌طبقه‌ای و حاشیه‌های تصمیم را تعیین نمی‌کند. ازاین‌رو شناخت درجات آزادی باقی‌مانده می‌تواند مقدمه‌ای برای بهبود هندسه پس از تعیین مراکز باشد.

\paragraph{خانوادهٔ پارامتری و تثبیت سیاست میانی.}
برخلاف تحلیل خانوادهٔ ناپارامتری گره‌محور آزاد در زیربخش~\ref{cmp:subsec:ch3-policy-geometry-map}، در اینجا خانوادهٔ پارامتری معرفی‌شده در رابطه~\eqref{cmp:eq:ch3-parametric-score-matrix} را در نظر می‌گیریم. در این خانواده، امتیازهای آغاز نمونه‌های مختلف می‌توانند متفاوت باشند، اما از توابعی با پارامترهای مشترک حاصل می‌شوند. فرض می‌کنیم پارامترهای آغاز از پارامترهای میانی جدا باشند:
\begin{equation}\label{cmp:eq:ch3-geometry-parametric-start-split}
\theta=(\theta_{\bot},\theta_{\mathrm{mid}}).
\end{equation}
این جداسازی به معنی استقلال سیاست آغاز نمونه‌ها از یکدیگر نیست؛ بلکه تنها امکان تغییر پارامترهای آغاز بدون تغییر سیاست میانی را فراهم می‌کند. در ادامه، داده، گراف عمودی، پارامترهای گشت و $\theta_{\mathrm{mid}}$ ثابت‌اند. دامنهٔ پارامترهای آغاز را در تحلیل موضعی، مجموعه‌ای باز در فضای اقلیدسی فرض می‌کنیم و توابع امتیاز را در همسایگی نقطهٔ مورد بررسی به‌طور پیوسته مشتق‌پذیر می‌گیریم. تمام امتیازها در این دامنه متناهی‌اند.

برای سیاست مشترک، احتمال‌های آغاز از رابطه~\eqref{cmp:eq:ch3-fixed-policy-softmax} با محدودسازی به ویژگی‌های موجود نمونه حاصل می‌شوند:
\begin{equation}\label{cmp:eq:ch3-geometry-parametric-start-probs}
p_{s,b}(\theta_\bot):=\pi(b\mid s,\bot;\theta_\bot)
=\frac{\exp(z_{\bot,b}(X_s;\theta_\bot))}
{\sum_{a\in A_{X_s}}\exp(z_{\bot,a}(X_s;\theta_\bot))},
\qquad b\in A_{X_s}.
\end{equation}
برای خانوادهٔ طبقه‌محور، در گشت آموزشی از عضو متناظر با طبقهٔ واقعی مبدأ استفاده می‌شود؛ در صورت استقلال بلوک‌های پارامترهای طبقه‌ها، تحلیل به هر بلوک نیز اعمال می‌شود. در صورت اشتراک پارامتر میان طبقه‌ها، همهٔ بلوک‌های بازنمایی و مراکز باید هم‌زمان بررسی شوند.

با ثابت بودن سیاست میانی، نقاط $u_{s,b}^{(H)}$ در رابطه~\eqref{cmp:eq:ch3-geometry-initial-feature-points} ثابت‌اند. اگر سطرهای ماتریس $U_s\in\mathbb R^{d_s\times K}$ را به‌ترتیب ویژگی‌های موجود، همین نقاط قرار دهیم و $p_s$ را بردار سطری احتمال‌های آغاز بگیریم، مطابق رابطه~\eqref{cmp:eq:ch3-geometry-origin-representation} داریم:
\begin{equation}\label{cmp:eq:ch3-geometry-parametric-start-representation}
\mathscr Y_s(\theta_\bot)=p_s(\theta_\bot)U_s,
\qquad
\mathcal V_c(\theta_\bot)=\frac{1}{|C_c|}\sum_{s\in C_c}p_s(\theta_\bot)U_s.
\end{equation}
هرچند بازنمایی نسبت به احتمال‌های آغاز خطی است، نگاشت پارامترها به احتمال‌ها عموماً غیرخطی است. به‌ویژه، اشتراک پارامترها و تفاوت مجموعه‌های $A_{X_s}$ مانع از آن می‌شوند که احتمال‌های آغاز نمونه‌ها را مستقل از یکدیگر تغییر دهیم. بنابراین جمع مینکوفسکی رابطه~\eqref{cmp:eq:ch3-geometry-center-minkowski}، در حالت سیاست میانی ثابت، برای این خانوادهٔ مقید یک کران بیرونی است و لزوماً فضای دقیق مراکز دست‌یافتنی نیست. محدب بودن فضای جواب‌های خانوادهٔ گره‌محور آزاد نیز به خانوادهٔ پارامتری تعمیم داده نمی‌شود.

\paragraph{نگاشت‌ها و مجموعه‌های هم‌ارزش.}
بردار حاصل از کنار هم قرار دادن احتمال‌های مجاز آغاز تمام نمونه‌های مرجع را با $\pi_\bot(\theta_\bot)$ نشان می‌دهیم؛ مؤلفه‌های ناموجود در آن وارد نمی‌شوند. منظور از سیاست مؤثر در این تحلیل، مقادیر این احتمال‌ها روی دادهٔ مرجع است، نه برابری تابع سیاست روی تمام نمونه‌های ممکن خارج از مرجع. با ثابت بودن سیاست میانی، زنجیرهٔ نگاشت‌ها به شکل زیر است:
\begin{equation}\label{cmp:eq:ch3-geometry-parametric-start-chain}
\theta_\bot\longmapsto\pi_\bot(\theta_\bot)
\longmapsto\mathscr Y(\theta_\bot)
\longmapsto\mathcal V(\theta_\bot).
\end{equation}
دو نگاشت آخر خطی‌اند؛ نگاشت مراکز از بازنمایی‌ها همان رابطه~\eqref{cmp:eq:ch3-class-centers-matrix} است.

فرض کنیم $\theta_\bot^\star$ پارامتر شاهد و $\mathscr Y^\star,\mathcal V^\star,\pi_\bot^\star$ خروجی‌های متناظر با آن باشند. مجموعه‌های پارامترهای هم‌سیاست، هم‌بازنمایی و هم‌مرکز را به‌ترتیب تعریف می‌کنیم:
\begin{equation}\label{cmp:eq:ch3-geometry-parametric-fiber-sets}
\begin{aligned}
\mathfrak F_\pi^\bot&=\{\theta_\bot:\pi_\bot(\theta_\bot)=\pi_\bot^\star\},\\
\mathfrak F_{\mathscr Y}^\bot&=\{\theta_\bot:\mathscr Y(\theta_\bot)=\mathscr Y^\star\},\\
\mathfrak F_{\mathcal V}^\bot&=\{\theta_\bot:\mathcal V(\theta_\bot)=\mathcal V^\star\}.
\end{aligned}
\end{equation}
وابستگی این مجموعه‌ها به نقطهٔ شاهد و پارامترهای ثابت میانی در نمادگذاری ضمنی است. مجموعهٔ آخر همتای پارامتری مجموعهٔ هم‌مرکز رابطه~\eqref{cmp:eq:ch3-fixed-center-fiber} با محدود کردن تغییرات به پارامترهای آغاز است. هر سه مجموعه ناتهی‌اند و
\begin{equation}\label{cmp:eq:ch3-geometry-parametric-fiber-inclusion}
\theta_\bot^\star\in\mathfrak F_\pi^\bot
\subseteq\mathfrak F_{\mathscr Y}^\bot\subseteq\mathfrak F_{\mathcal V}^\bot.
\end{equation}
این شمول سراسری است؛ اما هیچ حکم عمومی دربارهٔ محدب بودن، اتصال‌پذیری یا تک‌عضوی بودن این مجموعه‌ها از آن نتیجه نمی‌شود.

\paragraph{ژاکوبین نگاشت‌ها.}
گرادیان هر امتیاز نسبت به $\theta_\bot$ را بردار ستونی می‌گیریم و برای خروجی‌ها نیز قرارداد برداری ستونی را به کار می‌بریم. با مشتق‌گیری از رابطه~\eqref{cmp:eq:ch3-geometry-parametric-start-probs} داریم:
\begin{equation}\label{cmp:eq:ch3-geometry-parametric-start-prob-gradient}
\nabla_{\theta_\bot}p_{s,b}=p_{s,b}\left(\nabla_{\theta_\bot}z_{\bot,b}(X_s;\theta_\bot)-\sum_{a\in A_{X_s}}p_{s,a}\nabla_{\theta_\bot}z_{\bot,a}(X_s;\theta_\bot)\right).
\end{equation}
با جایگذاری در مشتق رابطه~\eqref{cmp:eq:ch3-geometry-parametric-start-representation} به دست می‌آید:
\begin{equation}\label{cmp:eq:ch3-geometry-parametric-start-Jsample}
\boxed{J_{\mathscr Y_s}=\sum_{b\in A_{X_s}}p_{s,b}
\bigl(u_{s,b}^{(H)}-\mathscr Y_s\bigr)^\top
\bigl(\nabla_{\theta_\bot}z_{\bot,b}(X_s;\theta_\bot)\bigr)^\top.}
\end{equation}
این ماتریس $K\times|\theta_\bot|$ است. حساسیت بازنمایی به پارامترهای آغاز هم به تفاوت نقاط هندسی ویژگی‌های آغازین و هم به حساسیت توابع امتیاز وابسته است. وابستگی پارامترها میان نمونه‌ها با به‌کارگیری همان ستون‌های پارامتر در ژاکوبین تمام نمونه‌ها حفظ می‌شود.

ژاکوبین‌های سه نگاشت خروجی نسبت به پارامترهای آغاز را با
\begin{equation}\label{cmp:eq:ch3-geometry-parametric-start-jacobians}
J_\pi=\frac{\partial\pi_\bot}{\partial\theta_\bot},\qquad
J_{\mathscr Y}=\frac{\partial\operatorname{vec}(\mathscr Y)}{\partial\theta_\bot},\qquad
J_{\mathcal V}=\frac{\partial\operatorname{vec}(\mathcal V)}{\partial\theta_\bot}
\end{equation}
نشان می‌دهیم. درایه‌های $\pi_\bot$ در این رابطه به شکل یک بردار ستونی مرتب شده‌اند. در هر قرارداد سازگار برداری‌سازی، نگاشت خطی از احتمال‌ها به بازنمایی‌ها و سپس به مراکز را می‌توان با ماتریس‌هایی نمایش داد؛ بنابراین طبق قاعدهٔ زنجیره‌ای، $J_{\mathscr Y}$ حاصل ضرب یک ماتریس ثابت در $J_\pi$ و $J_{\mathcal V}$ حاصل ضرب یک ماتریس ثابت در $J_{\mathscr Y}$ است. در نتیجه:
\begin{equation}\label{cmp:eq:ch3-geometry-parametric-kernel-inclusion}
\boxed{\ker J_\pi\subseteq\ker J_{\mathscr Y}\subseteq\ker J_{\mathcal V}.}
\end{equation}
با تعریف رتبه‌ها به‌صورت
\begin{equation}\label{cmp:eq:ch3-geometry-parametric-ranks}
r_\pi:=\operatorname{rank}J_\pi,\qquad
r_{\mathscr Y}:=\operatorname{rank}J_{\mathscr Y},\qquad
r_{\mathcal V}:=\operatorname{rank}J_{\mathcal V},
\end{equation}
از رابطهٔ شمول نتیجه می‌شود که
\begin{equation}\label{cmp:eq:ch3-geometry-parametric-rank-inclusion}
0\le r_{\mathcal V}\le r_{\mathscr Y}\le r_\pi\le|\theta_\bot|.
\end{equation}

\paragraph{تفکیک درجات آزادی مرتبهٔ اول.}
فضاهای پوچ فوق، جهت‌های تغییر پارامتر را مشخص می‌کنند که به‌ترتیب سیاست مؤثر، بازنمایی‌ها و مراکز را در مرتبهٔ اول ثابت نگه می‌دارند. برای تمایز افزونگی پارامترسازی از چندگانگی واقعی سیاست و هندسه، باید فضای پوچ کوچک‌تر را از فضای پوچ بزرگ‌تر خارج‌قسمت بگیریم. به‌طور مشخص، از قضیهٔ رتبه و پوچی داریم:
\begin{equation}\label{cmp:eq:ch3-geometry-parametric-freedom-decomposition}
\boxed{\begin{aligned}
\dim\ker J_\pi&=|\theta_\bot|-r_\pi,\\
\dim\bigl(\ker J_{\mathscr Y}/\ker J_\pi\bigr)&=r_\pi-r_{\mathscr Y},\\
\dim\bigl(\ker J_{\mathcal V}/\ker J_{\mathscr Y}\bigr)&=r_{\mathscr Y}-r_{\mathcal V}.
\end{aligned}}
\end{equation}
سطر نخست آزادی تغییر پارامتر بدون تغییر مرتبهٔ اول سیاست، سطر دوم آزادی مؤثر تغییر سیاست با حفظ بازنمایی‌ها، و سطر سوم آزادی تغییر بازنمایی‌ها با حفظ مراکز را بیان می‌کند. بعد فضای تغییرات مرتبهٔ اول مراکز نیز $r_{\mathcal V}$ است. مجموع این چهار مقدار برابر $|\theta_\bot|$ خواهد بود.

برای اثبات تعبیر سطر دوم، نگاشت $J_\pi$ را به $\ker J_{\mathscr Y}$ محدود می‌کنیم. هستهٔ این محدودسازی دقیقاً $\ker J_\pi$ است؛ بنابراین بعد تصویر برابر $\dim\ker J_{\mathscr Y}-\dim\ker J_\pi=r_\pi-r_{\mathscr Y}$ است. به همین ترتیب، محدودسازی $J_{\mathscr Y}$ به $\ker J_{\mathcal V}$ دارای هستهٔ $\ker J_{\mathscr Y}$ است و بعد تصویر آن برابر $r_{\mathscr Y}-r_{\mathcal V}$ خواهد بود. این استدلال در هر نقطهٔ مشتق‌پذیری برقرار است و به فرض رتبهٔ ثابت نیاز ندارد؛ اما در این مرحله صرفاً دربارهٔ تغییرات مرتبهٔ اول صحبت می‌کنیم.

\paragraph{کران‌های رتبه و وابستگی به داده.}
به دلیل نرمال بودن بردار احتمال هر نمونه، تصویر ژاکوبین سیاست در حاصل‌ضرب فضاهای مماس سیمپلکس‌های ویژگی‌ها قرار می‌گیرد. همچنین هر بازنمایی و هر مرکز در سیمپلکس برچسب‌ها قرار دارد. ازاین‌رو:
\begin{equation}\label{cmp:eq:ch3-geometry-parametric-rank-bounds}
\boxed{\begin{aligned}
r_\pi&\le\min\left\{|\theta_\bot|,\sum_{s=1}^{n}(d_s-1)\right\},\\
r_{\mathscr Y}&\le\min\{r_\pi,n(K-1)\},\\
r_{\mathcal V}&\le\min\{r_{\mathscr Y},K(K-1)\}.
\end{aligned}}
\end{equation}
حتی می‌توان کران دوم را با استفاده از بعد آفین پوش محدب هر نمونه در رابطه~\eqref{cmp:eq:ch3-geometry-sample-polytope} به $r_{\mathscr Y}\le\min\{r_\pi,\sum_s\dim\operatorname{aff}(\mathcal P_s)\}$ تنگ‌تر کرد. به همین ترتیب، جمع ابعاد آفین چندوجهی‌های مراکز گره‌محور کران بالایی برای رتبهٔ ژاکوبین مراکز پارامتری است.

برای دو نوع آزادی مؤثر، از روابط بالا نتیجه می‌شود:
\begin{equation}\label{cmp:eq:ch3-geometry-parametric-effective-bounds}
\begin{aligned}
\max\{0,r_\pi-n(K-1)\}&\le r_\pi-r_{\mathscr Y}\le r_\pi,\\
\max\{0,r_{\mathscr Y}-K(K-1)\}&\le r_{\mathscr Y}-r_{\mathcal V}\le r_{\mathscr Y}.
\end{aligned}
\end{equation}
از سوی دیگر، برای فضای پوچ ژاکوبین مراکز، کران‌های مستقل از شمار نمونه‌ها (با تثبیت تعداد پارامترها و طبقه‌ها) عبارت‌اند از:
\begin{equation}\label{cmp:eq:ch3-geometry-parametric-center-null-bounds}
\max\{0,|\theta_\bot|-K(K-1)\}\le\dim\ker J_{\mathcal V}\le|\theta_\bot|.
\end{equation}
استقلال این کران‌ها از $n$ به معنی استقلال رتبه‌ها یا بعد واقعی مجموعه‌ها از داده و تعداد نمونه‌ها نیست. افزونگی پارامترسازی نیز ممکن است تمام فضای پوچ را اشغال کند؛ بنابراین بزرگ بودن $|\theta_\bot|$ به‌تنهایی وجود آزادی مؤثر سیاست یا هندسه را تضمین نمی‌کند. بدون فرض‌های اضافی دربارهٔ توابع امتیاز و تنوع نقاط هندسی، کران پایین مثبت عمومی برای دو اختلاف رتبه به دست نمی‌آید.

\paragraph{از فضای پوچ تا مجموعه‌های دقیق هم‌ارزش.}
صفر شدن مشتق در یک جهت لزوماً به معنی وجود مسیر غیرثابتی از جواب‌های دقیق نیست. برای مثال، نگاشت حقیقی $x\mapsto x^2$ در صفر مشتق صفر دارد، ولی پیش‌تصویر صفر تنها یک عضو دارد. بنابراین برای عبور از تحلیل خطی موضعی به مجموعهٔ دقیق جواب‌ها، به شرطی اضافی نیاز است.

فرض کنیم رتبهٔ $J_{\mathscr Y}$ در یک همسایگی باز از $\theta_\bot^\star$ ثابت باشد. طبق قضیهٔ رتبهٔ ثابت، در یک همسایگی به‌اندازهٔ کافی کوچک، مجموعهٔ $\mathfrak F_{\mathscr Y}^\bot$ خمینه‌ای هموار از ردهٔ $C^1$ است، با
\begin{equation}\label{cmp:eq:ch3-geometry-parametric-Y-fiber-dimension}
\dim_{\mathrm{local}}\mathfrak F_{\mathscr Y}^\bot=|\theta_\bot|-r_{\mathscr Y},
\qquad
T_{\theta_\bot^\star}\mathfrak F_{\mathscr Y}^\bot=\ker J_{\mathscr Y}(\theta_\bot^\star).
\end{equation}
اگر در نقطهٔ شاهد $r_\pi>r_{\mathscr Y}$ باشد، بنا بر رابطه~\eqref{cmp:eq:ch3-geometry-parametric-freedom-decomposition} بردار $v$ وجود دارد که $J_{\mathscr Y}v=0$ ولی $J_\pi v\ne0$. چون $v$ در فضای مماس خمینه است، مسیری مشتق‌پذیر $\theta_\bot(t)$ بر روی آن وجود دارد که در $t=0$ از نقطهٔ شاهد می‌گذرد و مشتق آن $v$ است. در یک بازهٔ باز کوچک حول صفر داریم:
\begin{equation}\label{cmp:eq:ch3-geometry-parametric-exact-Y-path}
\mathscr Y(\theta_\bot(t))=\mathscr Y^\star,
\qquad
\left.\frac{\mathrm d\pi_\bot(\theta_\bot(t))}{\mathrm dt}\right|_{t=0}=J_\pi v\ne0.
\end{equation}
چون دست‌کم یک مؤلفهٔ سیاست مشتق ناصفر دارد و مشتق آن پیوسته است، با کوچک کردن بازه می‌توان آن مؤلفه را در طول مسیر اکیداً یکنواخت کرد. بنابراین در این بازه بی‌نهایت (در واقع ناشمارا) سیاست مؤثر متمایز با بازنمایی‌های دقیقاً یکسان وجود دارد. برای این نتیجه، ثابت بودن رتبهٔ $J_\pi$ لازم نیست؛ ثابت بودن رتبهٔ نگاشت بازنمایی و اختلاف رتبه در نقطهٔ شاهد کافی است.

با استدلالی مشابه، اگر رتبهٔ $J_{\mathcal V}$ در همسایگی نقطهٔ شاهد ثابت باشد، داریم:
\begin{equation}\label{cmp:eq:ch3-geometry-parametric-V-fiber-dimension}
\dim_{\mathrm{local}}\mathfrak F_{\mathcal V}^\bot=|\theta_\bot|-r_{\mathcal V},
\qquad
T_{\theta_\bot^\star}\mathfrak F_{\mathcal V}^\bot=\ker J_{\mathcal V}(\theta_\bot^\star).
\end{equation}
اگر در نقطهٔ شاهد $r_{\mathscr Y}>r_{\mathcal V}$ باشد، برداری $v$ با $J_{\mathcal V}v=0$ و $J_{\mathscr Y}v\ne0$ وجود دارد و روی خمینهٔ هم‌مرکز مسیری خواهیم داشت که
\begin{equation}\label{cmp:eq:ch3-geometry-parametric-exact-V-path}
\mathcal V(\theta_\bot(t))=\mathcal V^\star,
\qquad
\left.\frac{\mathrm d\operatorname{vec}(\mathscr Y(\theta_\bot(t)))}{\mathrm dt}\right|_{t=0}=J_{\mathscr Y}v\ne0.
\end{equation}
به این ترتیب، وجود مسیر پیوسته‌ای از آرایش‌های متمایز بازنمایی با مراکز دقیقاً یکسان اثبات می‌شود. این نتیجه به محدب بودن فضای سیاست‌ها، یکتایی پارامترسازی یا استقلال احتمال‌های آغاز نمونه‌ها وابسته نیست.

در صورت ثابت بودن رتبهٔ هر سه ژاکوبین در همسایگی، روابط اختلاف رتبه نه‌تنها بعد تصویر تغییرات مرتبهٔ اول، بلکه بعد موضعی تصویر سیاست روی خمینهٔ هم‌بازنمایی و تصویر بازنمایی روی خمینهٔ هم‌مرکز را نیز مشخص می‌کنند. دلیل آن است که ژاکوبین نگاشت محدودشده به هر خمینه در آن همسایگی دارای رتبهٔ ثابت، برابر اختلاف رتبه‌های متناظر است؛ در نتیجه قضیهٔ رتبهٔ ثابت بر نگاشت محدودشده نیز اعمال می‌شود. این نتیجه محلی است و وجود شاخه‌های دیگر یا چندگانگی سراسری را نفی نمی‌کند.

\paragraph{شرایط کافی چندگانگی و یکتایی موضعی.}
با استفاده از کران‌های رابطه~\eqref{cmp:eq:ch3-geometry-parametric-rank-bounds}، دو شرط کافی برای مثبت بودن اختلاف رتبه‌های مورد نظر به دست می‌آید:
\begin{equation}\label{cmp:eq:ch3-geometry-parametric-multiple-sufficient}
\begin{aligned}
r_\pi>n(K-1)&\ \Longrightarrow\ r_\pi>r_{\mathscr Y},\\
r_{\mathscr Y}>K(K-1)&\ \Longrightarrow\ r_{\mathscr Y}>r_{\mathcal V}.
\end{aligned}
\end{equation}
برای نتیجه‌گیری وجود مسیر دقیق، همچنان شرط رتبهٔ ثابت نگاشت بازنمایی یا مراکز، به‌ترتیب، لازم است. این شرط‌های کافی برای روش پیشنهادی به معنای ضرورت آن‌ها برای چندگانگی نیستند و تحققشان برای یک داده یا خانوادهٔ دلخواه تضمین نشده است.

در جهت مقابل، اگر در همسایگی نقطهٔ شاهد رتبهٔ نگاشت بازنمایی ثابت و $r_\pi=r_{\mathscr Y}$ باشد، مشتق سیاست روی خمینهٔ هم‌بازنمایی صفر است. با کوچک کردن همسایگی به یک وصلهٔ متصل از خمینه، سیاست روی آن ثابت خواهد بود. به همین ترتیب، اگر رتبهٔ نگاشت مراکز ثابت و $r_{\mathscr Y}=r_{\mathcal V}$ در همسایگی برقرار باشد، بازنمایی‌ها روی وصلهٔ محلی هم‌مرکز ثابت خواهند بود. بنابراین به‌ترتیب یکتایی موضعی سیاست مؤثر با تثبیت بازنمایی و یکتایی موضعی آرایش با تثبیت مرکز حاصل می‌شود. برابری رتبه‌ها فقط در یک نقطه برای این نتیجه کافی نیست. همچنین این یکتایی به پیش‌تصویرهای واقع در همسایگی پارامتر شاهد مربوط است و نه به تمام شاخه‌های ممکن در فضای پارامترها.

شرط کافی قوی‌تری برای یکتایی خود پارامترهای آغاز عبارت است از:
\begin{equation}\label{cmp:eq:ch3-geometry-parametric-full-rank-unique}
\operatorname{rank}J_{\mathcal V}(\theta_\bot^\star)=|\theta_\bot|.
\end{equation}
در این صورت می‌توان $|\theta_\bot|$ مؤلفهٔ خروجی را با ژاکوبین مربعی وارون‌پذیر انتخاب کرد؛ قضیهٔ تابع معکوس برای این زیرنگاشت، یک‌به‌یک بودن نگاشت کامل مراکز را در همسایگی نقطهٔ شاهد نتیجه می‌دهد. شرط لازم برای رتبهٔ ستونی کامل، $|\theta_\bot|\le K(K-1)$ است. در حضور ناوردایی پارامترسازی (مانند انتقال هم‌زمان همهٔ امتیازهای یک سطر)، باید ابتدا افزونگی را با یک قرارداد شناسایی مناسب حذف کرد؛ رتبهٔ کامل در فضای خام الزاماً امکان‌پذیر نیست. برای نگاشت بازنمایی و سیاست نیز استدلال رتبهٔ ستونی کامل به‌طور مشابه برقرار است.

\paragraph{دربارهٔ فرض رتبهٔ ثابت.}
اگر توابع امتیاز نسبت به پارامترها حقیقی‌تحلیلی باشند، با ثابت بودن ماسک‌ها و متناهی بودن امتیازها، نگاشت‌های سیاست، بازنمایی و مراکز نیز حقیقی‌تحلیلی‌اند. روی هر مؤلفهٔ باز و متصل از دامنه، برای هر ژاکوبین با بیشینهٔ رتبهٔ $r>0$، یک مینور مرتبهٔ $r$ وجود دارد که به‌عنوان تابع تحلیلی، همانی صفر نیست. مجموعهٔ صفر این مینور اندازهٔ لبگ صفر دارد و خارج از آن، رتبه برابر بیشینهٔ رتبه و در همسایگی نقطه ثابت است. برای $r=0$ نیز رتبه در کل مؤلفه ثابت است. در نتیجه، برای تعداد متناهی این نگاشت‌ها، شرط رتبهٔ ثابت به‌طور هم‌زمان تقریباً همه‌جا برقرار است. این استدلال به خاصیت شناخته‌شدهٔ مجموعهٔ صفر توابع حقیقی‌تحلیلی ناصفر متکی است؛ برای توابع صرفاً $C^1$ چنین نتیجهٔ تقریباً همه‌جایی در حالت عمومی ادعا نمی‌شود. مهم‌تر آنکه نقاط حاصل از بهینه‌سازی می‌توانند در مجموعهٔ استثنایی قرار گیرند؛ پس این مشاهده جای بررسی شرط رتبهٔ ثابت در نقطهٔ شاهد را نمی‌گیرد.

\paragraph{حالت خاص سیاست ثابت ماتریسی: یکتایی و چندگانگی سراسری.}
برای شناخت دقیق‌تر شرایط یکتایی، خانوادهٔ ثابت ماتریسی رابطه~\eqref{cmp:eq:ch3-fixed-policy-matrix} را با تثبیت سیاست میانی بررسی می‌کنیم. ابتدا حالت طبقه‌محور با سطر آغاز مستقل برای هر طبقه و بدون ویژگی ناموجود را در نظر می‌گیریم. تعداد امتیازهای آغاز هر طبقه $d$ است و به دلیل ناوردایی نرمال‌سازی با بیشینهٔ هموار نسبت به جمع یک ثابت، تعداد درجات آزادی مؤثر آن $d-1$ است. تعداد طبقه‌ها $K$ لزوماً برابر $d$ نیست.

چون همهٔ نمونه‌ها مجموعهٔ یکسانی از ویژگی‌ها دارند، بردار احتمال آغاز $p^{(c)}\in\operatorname{relint}(\Delta_d)$ در تمام نمونه‌های طبقهٔ $c$ یکسان است. با ترتیب ثابت ویژگی‌ها، تعریف می‌کنیم:
\begin{equation}\label{cmp:eq:ch3-geometry-fixed-start-average-U}
\overline U_c:=\frac{1}{|C_c|}\sum_{s\in C_c}U_s,\qquad
\mathcal V_c=p^{(c)}\overline U_c.
\end{equation}
در خانوادهٔ طبقه‌محور، $U_s$ با استفاده از سیاست میانی همان طبقه ساخته می‌شود. برای یک مرکز دست‌یافتنی $\mathcal V_c^\star$، مجموعهٔ احتمال‌های آغاز بازتولیدکنندهٔ آن برابر است با
\begin{equation}\label{cmp:eq:ch3-geometry-fixed-start-center-fiber}
\{p^{(c)}\in\operatorname{relint}(\Delta_d):p^{(c)}\overline U_c=\mathcal V_c^\star\}.
\end{equation}
این مجموعه محدب است. چون $\overline U_c\mathbf1_K=\mathbf1_d$، شرط $\delta p^{(c)}\overline U_c=0$ نرمال ماندن احتمال‌ها را نیز نتیجه می‌دهد. برای شاهد اکیداً مثبت، حرکت‌های کافی کوچک در تمام جهت‌های پوچ ممکن‌اند؛ بنابراین بعد آفین مجموعهٔ جواب‌ها دقیقاً برابر است با
\begin{equation}\label{cmp:eq:ch3-geometry-fixed-start-center-dim}
\boxed{d-\operatorname{rank}(\overline U_c).}
\end{equation}
از این رابطه سه نتیجه حاصل می‌شود: اگر $d>K$ باشد، به دلیل $\operatorname{rank}(\overline U_c)\le K$، چندگانگی احتمال‌های آغاز تضمین می‌شود؛ اگر $d\le K$ باشد، شرط لازم و کافی یکتایی برای هر شاهد اکیداً مثبت، رتبهٔ سطری کامل $\operatorname{rank}(\overline U_c)=d$ است. در حالت خاص $d=K$ و رتبهٔ کامل، احتمال آغاز یکتا به شکل
\begin{equation}\label{cmp:eq:ch3-geometry-fixed-start-inverse}
p^{(c)}=\mathcal V_c^\star\overline U_c^{-1}
\end{equation}
به دست می‌آید. این یکتایی سراسری و در فضای احتمال‌هاست، نه در فضای امتیازهای خام. برای هدفی که دست‌یافتنی بودن آن معلوم نیست، باید اعتبار بردار حاصل از وارون‌سازی نیز بررسی شود.

برای حفظ تمام بازنمایی‌های طبقه، ماتریس‌ها را به‌صورت افقی کنار هم قرار می‌دهیم: $\mathsf U_c=[U_s]_{s\in C_c}\in\mathbb R^{d\times (K|C_c|)}$. شرط حفظ تمام بازنمایی‌ها $\delta p^{(c)}\mathsf U_c=0$ است و بعد مجموعهٔ احتمال‌های هم‌بازنمایی برابر $d-\operatorname{rank}(\mathsf U_c)$ خواهد بود. از آنجا که $\overline U_c$ ترکیبی خطی از بلوک‌های $\mathsf U_c$ است، اختلاف $\operatorname{rank}(\mathsf U_c)-\operatorname{rank}(\overline U_c)$ دقیقاً بعد آفین آرایش‌های هم‌مرکز دست‌یافتنی در این حالت است. بنابراین حتی در صورت چندگانگی سیاست‌های هم‌مرکز، امکان تغییر هندسه به اختلاف این دو رتبه بستگی دارد.

اگر سیاست ثابت میان همهٔ طبقه‌ها مشترک باشد، یک بردار احتمال آغاز مشترک داریم؛ شرط حفظ همهٔ مراکز به‌صورت دستگاه خطی با ماتریس افقی $[\overline U_1,\ldots,\overline U_K]$ نوشته می‌شود. در این حالت نیز یکتایی معادل رتبهٔ سطری کامل این ماتریس (برای شاهد مثبت) است، اما شرط $d>K$ به‌تنهایی برای چندگانگی کافی نیست، زیرا قیدهای همهٔ مراکز بر همان بردار احتمال مشترک اعمال می‌شوند.

\paragraph{اثر ماسک در سیاست ثابت.}
اگر نمونه‌ها مجموعه‌های متفاوتی از ویژگی‌ها داشته باشند، حتی با امتیازهای ثابت مشترک، احتمال‌های نرمال‌شدهٔ آغاز میان نمونه‌ها یکسان نیستند. برای یک عضو ثابت خانواده، $p$ را بردار احتمال‌های مثبت نرمال‌شده روی اجتماع ویژگی‌های مورد استفادهٔ آن عضو می‌گیریم؛ بُعد این بردار را در این بند با $d'$ نشان می‌دهیم. در سیاست طبقه‌محور می‌توان آن را فقط روی اجتماع ویژگی‌های اعضای همان طبقه تعریف کرد. در این صورت
\begin{equation}\label{cmp:eq:ch3-geometry-fixed-masked-start}
\mathscr Y_s(p)=\frac{\sum_{b\in A_{X_s}}p_b u_{s,b}^{(H)}}{\sum_{b\in A_{X_s}}p_b}.
\end{equation}
برای حفظ بازنمایی معلوم $\mathscr Y_s^\star$، با ضرب در مخرج مثبت، رابطهٔ خطی زیر حاصل می‌شود:
\begin{equation}\label{cmp:eq:ch3-geometry-fixed-masked-representation-fiber}
\sum_{b\in A_{X_s}}p_b\bigl(u_{s,b}^{(H)}-\mathscr Y_s^\star\bigr)=0.
\end{equation}
با کنار هم قرار دادن این قیود برای همهٔ نمونه‌های مربوط، دستگاه $Mp^\top=0$ به دست می‌آید. با توجه به وجود شاهد مثبت و قید $p\mathbf1_{d'}=1$، مجموعهٔ جواب‌ها در فضای $p$ محدب است و بعد آفین آن برابر $d'-1-\operatorname{rank}(M)$ است. قید نرمال بودن از سطرهای $M$ مستقل است، زیرا یک شاهد در فضای پوچ $M$ وجود دارد که مجموع مؤلفه‌های آن یک است. توجه شود که در حضور ماسک‌ها، تغییر $p$ لزوماً به تغییر سیاست مؤثر روی دادهٔ مرجع نمی‌انجامد (برای مثال اگر گروه‌هایی از ویژگی‌ها هرگز در یک نمونه با هم ظاهر نشوند، مقیاس نسبی آن گروه‌ها ممکن است قابل شناسایی نباشد). بنابراین این بعد را نباید بدون بررسی افزونگی، بعد فضای سیاست‌های مؤثر دانست.

برای حفظ صرفاً مراکز، میانگین روابط~\eqref{cmp:eq:ch3-geometry-fixed-masked-start} شامل مجموع کسرهایی با مخرج‌های متفاوت است؛ در نتیجه استدلال خطی و محدب بودن مجموعهٔ هم‌بازنمایی به مجموعهٔ هم‌مرکز تعمیم نمی‌یابد. این حالت خاص، علت اهمیت ماسک‌ها و اشتراک پارامترها در تحلیل عمومی را روشن می‌کند.

\paragraph{جمع‌بندی و زمینهٔ بهینه‌سازی هندسی.}
نتایج بالا نشان می‌دهند که بسته به خانوادهٔ سیاست، داده و نقطهٔ شاهد، هم یکتایی و هم چندگانگی واقعی سیاست‌ها یا آرایش‌های هندسی ممکن‌اند. در حالت عمومی، شمول هسته‌ها و اختلاف رتبه‌ها اطلاعات مرتبهٔ اول را به دست می‌دهند؛ برای عبور به مسیرهای دقیق، شرط رتبهٔ ثابت یا استدلال دیگری دربارهٔ مجموعهٔ جواب‌ها ضروری است. در حالت خاص سیاست ثابت بدون ماسک، ساختار خطی امکان نتیجه‌گیری سراسری را نیز فراهم می‌آورد.

اگر تابع هدف آموزش صرفاً به بازنمایی‌ها و مراکز وابسته باشد و جمله‌ای وابسته به خود پارامترها یا سیاست نداشته باشد، روی مسیر هم‌بازنمایی مقدار آن ثابت خواهد ماند. اگر نقطهٔ شاهد کمینهٔ محلی چنین تابعی باشد، نقاط کافی نزدیک و هم‌ارزش در همان همسایگی نیز کمینهٔ محلی‌اند؛ زیرا هر یک همسایگی کوچکی درون همسایگی کمینهٔ اولیه دارند. البته وجود این نقاط، فرار از کمینهٔ محلی یا بهبود زیان را تضمین نمی‌کند. برای چنین نتیجه‌ای باید رفتار تابع هدف در جهت‌های دیگر، از جمله جهت تغییر سیاست میانی، جداگانه بررسی شود.

از سوی دیگر، وجود مسیر هم‌مرکز با بازنمایی‌های متغیر، صرفاً امکان تغییر هندسه را اثبات می‌کند. برای معیار هندسی $f$، با قرارداد مطلوب بودن مقدار کمتر، می‌توان مسئلهٔ زیر را مطرح کرد:
\begin{equation}\label{cmp:eq:ch3-geometry-parametric-center-constrained-objective}
\inf_{\theta_\bot\in\mathfrak F_{\mathcal V}^\bot}
 f\bigl(\mathscr Y(\theta_\bot),\mathcal V^\star\bigr).
\end{equation}
وجود شاهد تضمین می‌کند این اینفیمم از مقدار هدف در نقطهٔ شاهد بیشتر نیست؛ اما بهبود اکید، تحقق اینفیمم و محدب بودن مسئله در پارامترهای مشترک به‌طور کلی تضمین نمی‌شوند. در ادامه باید بررسی کنیم برای معیارهای تراکم و فاصله از مجموعه‌های جداکننده، چه شرایطی امکان بهبود را فراهم می‌کنند و چگونه می‌توان از این شرایط در اصلاح سیاست استفاده کرد.

\paragraph{بهبود تراکم با حفظ دقیق مراکز.}
اکنون مسئلهٔ رابطه~\eqref{cmp:eq:ch3-geometry-parametric-center-constrained-objective} را برای معیارهای تراکم مرکزی معرفی‌شده در رابطه‌های~\eqref{cmp:eq:ch3-compactness-general} و~\eqref{cmp:eq:ch3-compactness-general-weighted} بررسی می‌کنیم. در این تحلیل $\psi\in\{\mathrm{KL},\mathrm{JS}\}$ است؛ منظور از سنجهٔ دوم خود واگرایی جنسن--شنون، نه جذر آن، و جهت KL از بازنمایی به مرکز است. وزن‌های طبقه‌ها ثابت و مثبت‌اند. در تمام روابط این بند، تابع هدف حتی خارج از مجموعهٔ هم‌مرکز نسبت به مراکز شاهد ثابت تعریف می‌شود:
\begin{equation}\label{cmp:eq:ch3-geometry-parametric-fixed-center-R}
R_\psi(\mathscr Y;\mathcal V^\star)
=\sum_{c=1}^K\frac{\omega_c}{|C_c|}
\sum_{s\in C_c}d_\psi(\mathscr Y_s,\mathcal V_c^\star).
\end{equation}
روی مجموعهٔ هم‌مرکز، این تابع دقیقاً همان معیار تراکم مرکزی فصل است. مشتق آن در جهت مماس نیز با مشتق معیار با مراکز متغیر یکسان است، زیرا مشتق مراکز در چنین جهتی صفر است.

برای هر نمونه، تکیه‌گاه ثابت $S_s^{\rm lab}:=\bigcup_{b\in A_{X_s}}\operatorname{supp}(u_{s,b}^{(H)})$ را در نظر می‌گیریم. در این بند، این نماد مجموعه‌ای از اندیس‌های برچسب است و با فضای حالت کاهش‌یافتهٔ $\mathcal S_s$ تفاوت دارد. به دلیل مثبت بودن احتمال‌های آغاز با امتیازهای متناهی، بازنمایی هر نمونه روی این تکیه‌گاه اکیداً مثبت و در خارج از آن صفر است. مرکز شاهد طبقه روی اجتماع تکیه‌گاه‌های اعضای آن طبقه مثبت است. ازاین‌رو مشتق‌گیری فقط نسبت به مؤلفه‌های $u\in S_s^{\rm lab}$ انجام می‌شود؛ مؤلفه‌های ساختاراً صفر در بردار گرادیان مورد استفاده در جهت‌های دست‌یافتنی، صفر قرار داده می‌شوند. این قرارداد به معنی مشتق‌پذیری واگرایی نسبت به خروج از یک وجه سیمپلکس نیست.

برای اختصار در فرمول مشتق‌ها، $x=\mathscr Y_s$ و $q=\mathcal V_c^\star$ را می‌گیریم، که $s\in C_c$. داریم:
\begin{equation}\label{cmp:eq:ch3-geometry-R-coordinate-gradient}
\begin{aligned}
\frac{\partial D_{\rm KL}(x\|q)}{\partial x_u}
&=1+\log\frac{x_u}{q_u},\\
\frac{\partial\operatorname{JSD}(x,q)}{\partial x_u}
&=\frac12\log\frac{2x_u}{x_u+q_u}.
\end{aligned}
\end{equation}
مشتق‌های متقاطع صفر و درایه‌های قطری هسین برابرند با
\begin{equation}\label{cmp:eq:ch3-geometry-R-coordinate-hessian}
\begin{aligned}
\frac{\partial^2 D_{\rm KL}(x\|q)}{\partial x_u^2}
&=\frac1{x_u},\\
\frac{\partial^2\operatorname{JSD}(x,q)}{\partial x_u^2}
&=\frac{q_u}{2x_u(x_u+q_u)}.
\end{aligned}
\end{equation}
هر دو در مؤلفه‌های فعال مثبت‌اند. در نتیجه این دو واگرایی نسبت به بازنمایی، با مرکز ثابت، روی وجه مجاز اکیداً محدب‌اند. این خاصیت با امتداد پیوستهٔ جمله‌های آنتروپی در صفر به وجه بسته نیز برقرار است: هر جملهٔ تک‌متغیره روی بازهٔ مثبت اکیداً محدب و در صفر پیوسته است و هر دو نقطهٔ متمایز دست‌کم در یک مؤلفه تفاوت دارند. بنابراین $R_\psi$ با وزن‌های مثبت روی حاصل‌ضرب وجوه مجاز اکیداً محدب است. این نتیجه برای جذر JSD یا وزن صفر طبقه به همان صورت ادعا نمی‌شود و نیز محدب بودن ترکیب آن با نگاشت پارامتری سیاست را نتیجه نمی‌دهد.

\paragraph{شرط لازم و کافی وجود جهت کاهشی مرتبهٔ اول.}
تمام ژاکوبین‌ها در نقطهٔ شاهد $\theta_\bot^\star$ ارزیابی می‌شوند، مگر آنکه خلاف آن تصریح شود. گرادیان ستونی معیار نسبت به بازنمایی‌های برداری‌شده را با $g_\psi$ و گرادیان آن نسبت به پارامترهای آغاز را با $g_\theta=J_{\mathscr Y}^{\top}g_\psi$ نشان می‌دهیم. در این قرارداد، فرافکن متعامد روی فضای پوچ ژاکوبین مراکز، همتای فرافکن رابطه~\eqref{cmp:eq:ch3-fixed-center-projection}، برابر است با
\begin{equation}\label{cmp:eq:ch3-geometry-parametric-center-projector}
\Pi_{\mathcal V}
=I_{|\theta_\bot|}-J_{\mathcal V}^{\dagger}J_{\mathcal V}.
\end{equation}
در اینجا $\dagger$ وارون مور--پنروز است و برداری‌سازی مراکز با ژاکوبین قبلی سازگار است. جهت $v=-\Pi_{\mathcal V}g_\theta$ دارای خواص زیر است:
\begin{equation}\label{cmp:eq:ch3-geometry-parametric-R-descent}
J_{\mathcal V}v=0,\qquad
g_\theta^\top v=-\|\Pi_{\mathcal V}g_\theta\|_2^2.
\end{equation}
به علاوه، برای هر $v\in\ker J_{\mathcal V}$ داریم $g_\theta^\top v=(\Pi_{\mathcal V}g_\theta)^\top v$. بنابراین
\begin{equation}\label{cmp:eq:ch3-geometry-parametric-R-descent-iff}
\exists v\in\ker J_{\mathcal V}:\ g_\theta^\top v<0
\quad\Longleftrightarrow\quad
\Pi_{\mathcal V}g_\theta\ne0.
\end{equation}
جهت فرافکنی‌شده در صورت ناصفر بودن، پس از نرمال‌سازی، بیشترین کاهش مشتق جهتی را میان بردارهای واحد فضای پوچ فراهم می‌کند؛ این تعبیر به متریک اقلیدسی انتخاب‌شده برای پارامترها وابسته است.

اگر علاوه بر این، رتبهٔ $J_{\mathcal V}$ در همسایگی شاهد ثابت باشد، طبق رابطه~\eqref{cmp:eq:ch3-geometry-parametric-V-fiber-dimension}، جهت $v$ مماس بر خمینهٔ دقیق هم‌مرکز است. پس مسیر $C^1$ مانند $\theta_\bot(t)$ روی این خمینه وجود دارد که $\theta_\bot(0)=\theta_\bot^\star$ و $\theta_\bot'(0)=v$. در نتیجه برای $t>0$ به‌اندازهٔ کافی کوچک، به شرط ناصفر بودن گرادیان فرافکنی‌شده،
\begin{equation}\label{cmp:eq:ch3-geometry-parametric-exact-R-improvement}
\mathcal V(\theta_\bot(t))=\mathcal V^\star,\qquad
R_\psi(\mathscr Y(\theta_\bot(t));\mathcal V^\star)
<R_\psi(\mathscr Y^\star;\mathcal V^\star).
\end{equation}
این نتیجه دربارهٔ وجود مسیر دقیق است و به معنی حفظ دقیق مرکز با یک گام مستقیم $t v$ نیست. اگر گرادیان فرافکنی‌شده صفر باشد، فقط نبود جهت کاهشی مرتبهٔ اول نتیجه می‌شود؛ بهینگی موضعی یا سراسری در حالت نامحدب عمومی از آن نتیجه نمی‌شود. جهت‌های کاهشی مرتبهٔ بالاتر نیز نفی نمی‌شوند.

\paragraph{جهت‌های انقباض کلی و زوجی.}
عملگر خطی میانگین‌گیری طبقه‌ای را در قرارداد برداری‌سازی سازگار با ژاکوبین‌ها با $\mathsf A$ نشان می‌دهیم؛ بنابراین $J_{\mathcal V}=\mathsf A J_{\mathscr Y}$. فضای تغییرات هندسی دست‌یافتنی مرتبهٔ اول با حفظ مراکز برابر است با
\begin{equation}\label{cmp:eq:ch3-geometry-parametric-geometric-tangent}
\mathcal T_{\mathscr Y}
=J_{\mathscr Y}(\ker J_{\mathcal V})
=\operatorname{Im}J_{\mathscr Y}\cap\ker\mathsf A,
\qquad
\dim\mathcal T_{\mathscr Y}=r_{\mathscr Y}-r_{\mathcal V}.
\end{equation}
برابری دو توصیف از $J_{\mathcal V}=\mathsf A J_{\mathscr Y}$ و بعد از محدودسازی ژاکوبین و قضیهٔ رتبه و پوچی حاصل می‌شود. به‌طور معادل، وجود جهت کاهشی مرتبهٔ اول به ناصفر بودن فرافکنی $g_\psi$ بر $\mathcal T_{\mathscr Y}$ بستگی دارد.

در فضای محیط، جهت انقباض کلی با مؤلفه‌های $h_s=\mathcal V_c^\star-\mathscr Y_s^\star$ برای $s\in C_c$ را در نظر می‌گیریم. چون $\sum_{s\in C_c}h_s=0$، مسیر مستقیم فرضی $\mathscr Y_s(t)=(1-t)\mathscr Y_s^\star+t\mathcal V_c^\star$ برای $0\le t\le1$ مرکز را ثابت نگه می‌دارد. با محدب بودن واگرایی و صفر بودن ناهمانندی مرکز با خودش:
\begin{equation}\label{cmp:eq:ch3-geometry-parametric-R-contraction}
R_\psi(\mathscr Y(t);\mathcal V^\star)
\le(1-t)R_\psi(\mathscr Y^\star;\mathcal V^\star).
\end{equation}
این مسیر در سیمپلکس معتبر است، اما لزوماً در فضای بازنمایی‌های دست‌یافتنی قرار ندارد. به‌ویژه ممکن است جهت $h_s$ در مؤلفه‌ای که برای نمونه ساختاراً صفر است، ناصفر باشد. اگر دستگاه
\begin{equation}\label{cmp:eq:ch3-geometry-parametric-R-contraction-lift}
J_{\mathscr Y}v=\operatorname{vec}(h)
\end{equation}
جواب داشته باشد، آن جهت با تکیه‌گاه‌های ثابت سازگار است، $J_{\mathcal V}v=0$ و با مشتق گرفتن از کران~\eqref{cmp:eq:ch3-geometry-parametric-R-contraction} در صفر داریم $g_\theta^\top v\le-R_\psi(\mathscr Y^\star;\mathcal V^\star)$. بنابراین برای تراکم غیرصفر، جهت کاهشی وجود دارد و تحت شرط رتبهٔ ثابت مراکز، بهبود دقیق رابطه~\eqref{cmp:eq:ch3-geometry-parametric-exact-R-improvement} حاصل می‌شود.

حتی دست‌یافتنی بودن یک جهت انقباض زوجی کافی است. برای دو عضو $i,j$ از یک طبقه با بازنمایی‌های متمایز، بردار $h$ را با $h_i=\mathscr Y_j^\star-\mathscr Y_i^\star$، $h_j=-h_i$ و صفر برای سایر اعضا تعریف می‌کنیم. این جهت میانگین طبقه را ثابت نگه می‌دارد. اگر در $\mathcal T_{\mathscr Y}$ قرار گیرد، تفاوت دو بازنمایی در تکیه‌گاه‌های مجاز هر دو نمونه قرار دارد. با اکیداً محدب بودن واگرایی با مرکز ثابت و یکنوایی اکید گرادیان روی تکیه‌گاه مشترک، برای گرادیان‌های ستونی $g_i,g_j$ واگرایی‌های دو نمونه داریم:
\begin{equation}\label{cmp:eq:ch3-geometry-parametric-R-pair-contraction}
g_\psi^\top\operatorname{vec}(h)
=-\frac{\omega_c}{|C_c|}
(g_i-g_j)^\top(\mathscr Y_i^\star-\mathscr Y_j^\star)^\top<0.
\end{equation}
در اینجا بردارهای $\mathscr Y_s$ سطری‌اند و $h$ به قرارداد برداری‌سازی قبلی تبدیل می‌شود. در صورت نبود تکیه‌گاه مشترک مناسب، انقباض فرضی را نباید جهت مجاز پارامتری فرض کرد.

وجود فضای غیرصفر $\mathcal T_{\mathscr Y}$ به تنهایی کاهش اکید تراکم را تضمین نمی‌کند. برای نمونهٔ هندسی در یک طبقه با سه بازنمایی دو مؤلفه‌ای $(0.3+t,0.7-t)$، $(0.3-t,0.7+t)$ و $(0.9,0.1)$، مرکز $(0.5,0.5)$ مستقل از $t$ است. در $t=0$ دو بازنمایی نخست منطبق‌اند و به دلیل تقارن و محدب بودن اکید، معیار تراکم KL یا JSD برای هر $t\ne0$ مجاز بیشتر از مقدار آن در صفر است. بنابراین با وجود تراکم غیرصفر و یک جهت هندسی آزاد، آرایش اولیه در این خانواده بهترین است. این مثال صرفاً هندسی است و تحقق آن توسط یک گراف مشخص IMC در اینجا ادعا نمی‌شود.

\paragraph{ظرفیت هندسهٔ پایه و محدودیت خانوادهٔ پارامتری.}
از چندوجهی $\mathcal P_s$ در رابطه~\eqref{cmp:eq:ch3-geometry-sample-polytope}، برای یک ویژگی دلخواه $b_0\in A_{X_s}$، فضای خطی جهت‌های بالقوه را تعریف می‌کنیم:
\begin{equation}\label{cmp:eq:ch3-geometry-parametric-base-direction-space}
E_s=\operatorname{span}\{u_{s,b}^{(H)}-u_{s,b_0}^{(H)}:b\in A_{X_s}\},
\quad r_s=\dim E_s=\dim\operatorname{aff}(\mathcal P_s).
\end{equation}
انتخاب $b_0$ اثری بر فضای حاصل ندارد. داریم $r_s\le\min\{d_s-1,K-1\}$. چون بازنمایی نمونه ترکیب نرمال‌شده‌ای از نقاط ثابت است، مشتق آن در $E_s$ قرار دارد. در نتیجه
\begin{equation}\label{cmp:eq:ch3-geometry-parametric-base-rank-bound}
\operatorname{Im}J_{\mathscr Y}\subseteq\prod_{s=1}^n E_s,\qquad
r_{\mathscr Y}\le\min\{r_\pi,\sum_s r_s\}.
\end{equation}
با استفاده از جمع مینکوفسکی رابطه~\eqref{cmp:eq:ch3-geometry-center-minkowski}، بعد آفین فضای مراکز هر طبقه در خانوادهٔ گره‌محور آزاد برابر است با $r_c:=\dim(\sum_{s\in C_c}E_s)$. چون عملگر میانگین‌گیری در هر طبقه تصویر $\prod_{s\in C_c}E_s$ را روی همین مجموع خطی می‌فرستد، داریم $r_{\mathcal V}\le\min\{r_{\mathscr Y},\sum_c r_c\}$.

فضای بالقوهٔ جهت‌های حافظ مراکز را تعریف می‌کنیم:
\begin{equation}\label{cmp:eq:ch3-geometry-parametric-base-center-direction-space}
\mathcal E=\left(\prod_{s=1}^nE_s\right)\cap\ker\mathsf A.
\end{equation}
بر اساس قضیهٔ رتبه و پوچی برای محدودسازی $\mathsf A$ به حاصل‌ضرب فضاها،
\begin{equation}\label{cmp:eq:ch3-geometry-parametric-base-center-direction-dim}
\dim\mathcal E=\sum_{s=1}^n r_s-\sum_{c=1}^K r_c,\qquad
\mathcal T_{\mathscr Y}\subseteq\mathcal E.
\end{equation}
ازاین‌رو اختلاف ابعاد
\begin{equation}\label{cmp:eq:ch3-geometry-parametric-capacity-gap}
\dim\mathcal E-\dim\mathcal T_{\mathscr Y}
=\sum_s r_s-\sum_c r_c-r_{\mathscr Y}+r_{\mathcal V}\ge0
\end{equation}
شمار جهت‌های بالقوهٔ حافظ مراکز را که به‌صورت مستقل در مرتبهٔ اول توسط خانوادهٔ پارامتری قابل ایجاد نیستند، مشخص می‌کند. این مقدار یک اختلاف بعد است، نه فاصلهٔ متریکی یا کران خطای تراکم. تمامی رتبه‌ها و فضاهای بالا به داده و سیاست میانی ثابت وابسته‌اند.

اگر شرط
\begin{equation}\label{cmp:eq:ch3-geometry-parametric-capacity-saturation}
r_{\mathscr Y}-r_{\mathcal V}=\sum_s r_s-\sum_c r_c
\end{equation}
برقرار باشد، به دلیل شمول و برابری ابعاد، $\mathcal T_{\mathscr Y}=\mathcal E$. این شرط را در این تحلیل اشباع ظرفیت هندسی حافظ مراکز می‌نامیم. اشباع فقط به فضای جهت‌های حافظ مراکز مربوط است و لزوماً مستلزم بیشینه بودن رتبهٔ ژاکوبین بازنمایی در همهٔ جهت‌ها نیست. شرط فوق قابل آزمون است؛ تحقق آن برای خانواده‌های پارامتری دلخواه یا نقاط حاصل از آموزش تضمین نشده است.

\paragraph{بهینگی هندسی در صورت اشباع ظرفیت.}
مجموعهٔ زیر را در قرارداد بستهٔ خانوادهٔ گره‌محور آزاد، با همان سیاست میانی، تعریف می‌کنیم:
\begin{equation}\label{cmp:eq:ch3-geometry-parametric-relaxed-center-arrangements}
\mathcal A(\mathcal V^\star)=
\{\mathscr Y:\mathscr Y_s\in\mathcal P_s\ \forall s,\ 
\mathsf A\operatorname{vec}(\mathscr Y)=\operatorname{vec}(\mathcal V^\star)\}.
\end{equation}
این مجموعه ناتهی، فشرده و محدب است. از آنجا که شاهد از احتمال‌های مثبت حاصل شده، $\mathscr Y_s^\star\in\operatorname{relint}(\mathcal P_s)$ برای همهٔ نمونه‌ها و بنابراین شاهد در درون نسبی مجموعهٔ مقید قرار دارد. به‌طور دقیق، برای هر $h\in\mathcal E$، با انتخاب $\varepsilon>0$ کافی کوچک، $\mathscr Y^\star\pm\varepsilon h$ در همین مجموعه قرار می‌گیرند (با قرارداد سازگار برداری‌سازی). در نتیجه $\mathcal E$ فضای جهت‌های آفین این مجموعه است.

با مثبت بودن همهٔ وزن‌های طبقه‌ها، پیوستگی و محدب بودن اکید معیار~\eqref{cmp:eq:ch3-geometry-parametric-fixed-center-R} روی این مجموعه، وجود و یکتایی آرایش
\begin{equation}\label{cmp:eq:ch3-geometry-parametric-relaxed-R-optimum}
\widehat{\mathscr Y}
=\arg\min_{\mathscr Y\in\mathcal A(\mathcal V^\star)}
R_\psi(\mathscr Y;\mathcal V^\star)
\end{equation}
را نتیجه می‌دهد. این کمینه‌ساز در فضای رهاسازی‌شده تعریف شده است و لزوماً با امتیازهای متناهی یا خانوادهٔ پارامتری مقید دست‌یافتنی نیست. به‌ویژه یکتایی آرایش به معنی یکتایی سیاست تولیدکنندهٔ آن نیست.

فرض کنیم شرط اشباع~\eqref{cmp:eq:ch3-geometry-parametric-capacity-saturation} در نقطهٔ شاهد برقرار باشد. اگر $\mathscr Y^\star\ne\widehat{\mathscr Y}$، بردار $h=\operatorname{vec}(\widehat{\mathscr Y}-\mathscr Y^\star)$ عضوی از $\mathcal E=\mathcal T_{\mathscr Y}$ خواهد بود. بنابراین $v\in\ker J_{\mathcal V}$ وجود دارد که $J_{\mathscr Y}v=h$. با نامساوی مرتبهٔ اول تابع محدب، حتی اگر کمینه‌ساز روی مرز مجموعهٔ بسته واقع شده باشد،
\begin{equation}\label{cmp:eq:ch3-geometry-parametric-saturated-descent}
g_\theta^\top v=g_\psi^\top h
\le R_\psi(\widehat{\mathscr Y};\mathcal V^\star)
-R_\psi(\mathscr Y^\star;\mathcal V^\star)<0.
\end{equation}
پس گرادیان فرافکنی‌شده ناصفر است. برعکس، اگر $\mathscr Y^\star=\widehat{\mathscr Y}$، چون درون نسبی مجموعهٔ مقید است، مشتق معیار در هر دو جهت $\pm h$ از فضای $\mathcal E$ باید صفر باشد. در نتیجه $g_\psi$ بر $\mathcal T_{\mathscr Y}$ عمود است و $g_\theta$ بر $\ker J_{\mathcal V}$ عمود خواهد بود. بنابراین
\begin{equation}\label{cmp:eq:ch3-geometry-parametric-saturated-stationarity}
\boxed{\Pi_{\mathcal V}g_\theta=0
\quad\Longleftrightarrow\quad
\mathscr Y^\star=\widehat{\mathscr Y}}
\end{equation}
تحت شرط اشباع و فرض‌های مشتق‌پذیری و وزن‌های مثبت برقرار است. این هم‌ارزی مرتبهٔ اول به‌خودی‌خود به ثابت بودن رتبه در همسایگی نیاز ندارد. برای نتیجه‌گیری وجود مسیر دقیق کاهشی در صورت نابرابری دو آرایش، شرط رتبهٔ ثابت ژاکوبین مراکز همچنان شرط کافی مورد استفاده است.

در نتیجه در هر نقطهٔ منظم از مجموعهٔ هم‌مرکز که ظرفیت هندسی اشباع شده است، کمینهٔ محلی نامطلوب نسبت به کمینهٔ مسئلهٔ محدب رهاسازی‌شده نمی‌تواند وجود داشته باشد. این حکم صرفاً دربارهٔ معیار منتخب با سیاست میانی و مراکز ثابت است؛ دربارهٔ کمینه‌های محلی کل مسئلهٔ آموزش یا نقاط تکین حکمی صادر نمی‌کند. اگر کمینهٔ رهاسازی‌شده دست‌یافتنی نباشد، هم‌ارزی بالا به معنای تضمین تحقق آن توسط سیاست پارامتری نیست.

\paragraph{حالت خاص سیاست ثابت بدون ماسک و نتیجهٔ کاربردی.}
در خانوادهٔ ثابت ماتریسی بدون تفاوت ماسک‌ها، با تثبیت سیاست میانی، بازنمایی‌ها نسبت به بردار احتمال آغاز مشترک (یا بردارهای مستقل طبقه‌ای) خطی‌اند و قیود حفظ مراکز خطی هستند. بنابراین کمینه‌سازی معیار~\eqref{cmp:eq:ch3-geometry-parametric-fixed-center-R} در فضای همین احتمال‌های آغاز یک مسئلهٔ محدب است و هر کمینهٔ محلی آن، کمینهٔ سراسری همان زیرمسئله است. در قرارداد بسته به دلیل فشردگی مجموعهٔ مجاز و پیوستگی هدف، کمینه محقق می‌شود؛ در قرارداد احتمال‌های مثبت ممکن است فقط اینفیمم وجود داشته باشد. در فضای بازنمایی‌های مجاز همین خانواده، آرایش بهینه در صورت تحقق یکتاست، اما سیاست لزوماً یکتا نیست. این نتیجه به معنی دست‌یابی به کمینهٔ خانوادهٔ آزاد گره‌محور نیست. با ماسک‌های متفاوت یا پارامترهای عمومی مشترک، این استدلال محدب بودن بدون فرض اضافی قابل تعمیم نیست.

به عنوان زمینهٔ پژوهشی برای ادامه، می‌توان آموزش نامحدب را با بهبود تراکم در زیرمسئلهٔ هم‌مرکز ترکیب کرد: ابتدا سیاست میانی و مراکز حاصل را ثابت نگه داریم، سپس سیاست آغاز را با قید مراکز بهبود دهیم و در پایان دوباره امکان تغییر سیاست میانی را فراهم کنیم. در خانوادهٔ عمومی، فرافکنی گرادیان فقط جهت مماس می‌دهد و یک روند محاسباتی حفظ مراکز به اصلاح غیرخطی یا کنترل باقیماندهٔ قیود نیاز دارد؛ در خانوادهٔ ثابت بدون ماسک، حل مستقیم مسئلهٔ محدب در فضای احتمال‌ها امکان‌پذیر است. بهبود تراکم با حفظ مراکز، بهبود قطعی زیان آموزشی، دقت برون‌نمونه‌ای یا خروج از یک کمینهٔ محلی واقعی تابع هدف اولیه را تضمین نمی‌کند. بررسی فاصله از مجموعه‌های جداکننده و حاشیه‌ها با مراکز ثابت، و طراحی روند اصلاح سیاست، موضوع ادامهٔ این زیربخش است.

% REVIEW NOTES (not part of thesis prose):
% 1. Check global notation and merge overlapping material from legacy section 3.10 only after editorial review.
% 2. Real-analytic generic-rank observation is conditional; add an appropriate reference.
% 3. Separate probability-space convexity from parameter-space and general masked-family results.
% 4. R analysis concerns forward KL and JSD, positive class weights, fixed active supports and finite scores. Review boundary and zero-weight cases separately; do not generalize to sqrt(JSD).
% 5. Constant rank is sufficient for lifting first-order directions to exact local paths, not asserted everywhere or necessary for every possible path.
% 6. Saturation criterion compares dimensions of center-preserving directions, not global attainability; the relaxed optimizer may lie on a boundary or be unattainable by the parametric family.
% 7. The three-sample counterexample is geometric only. No IMC graph witness has been claimed.
% 8. Possible empirical follow-up: section 4.14 noise experiments, fixed class-specific no-mask start probabilities, d=3+k and K=3 imply at least k center-preserving probability degrees of freedom per class; shared start lower bound max(0,k-4). Effective geometry changes require rank differences. These bounds do not establish the cause of observed optimization dynamics.
% 9. Next: fixed-center margins / signed distances, KL versus JS geometry; then exact correction, alternating optimization, and limitations on escape from local minima.
% 10. Retain all calculations for review; no migration to appendix or main chapter until editorial decision.
